\documentclass[preview]{standalone}

\usepackage{amsmath}
\usepackage{amssymb}
\usepackage{stellar}
\usepackage{definitions}

\begin{document}

\id{psicologia-della-gestalt}
\genpage

\section{La psicologia della Gestalt}

\begin{snippetdefinition}{gestalt-definition}{Gestalt}
    Il termine \emph{Gestalt}, cioè forma, indica una configurazione unitaria
    e organizzata che differisce dalla semplice somma delle sue singole parti.
\end{snippetdefinition}

\begin{snippet}{nascita-gestalt}
    La psicologia della Gestalt nasce nel 1912, quando \emph{Max Wertheimer}
    pubblica il lavoro sul movimento apparente, o fenomeno phi. Oltre a
    Wertheimer, i principali esponenti della Gestalt sono \emph{Koffka} e
    \emph{Köhler}.
\end{snippet}

\begin{snippetdefinition}{fenomeno-phi-definition}{Fenomeno phi}
    Il \emph{fenomeno phi} è un fenomeno percettivo per cui l'osservatore
    percepisce un unico stimolo in movimento a partire da stimoli luminosi
    presentati in successione.
\end{snippetdefinition}

\begin{snippetexample}{movimento-apparente-example}{Movimento apparente}
    Il movimento apparente si ottiene proiettando su uno sfondo scuro uno
    stimolo luminoso e, dopo un intervallo molto breve, un secondo stimolo
    luminoso vicino alla posizione del primo. L'osservatore non percepisce due
    stimoli separati, ma un unico stimolo in movimento. Questo fenomeno mostra
    che la percezione non coincide con la semplice somma degli stimoli fisici.
\end{snippetexample}

\begin{snippet}{gestalt-principi}
    La Gestalt rifiuta l'elementarismo e concentra l'attenzione su struttura e
    organizzazione. L'analisi deve procedere dall'intero alle sue parti
    costituenti. Costruire una teoria della percezione significa individuare le
    regole dell'interazione tra le parti del campo fenomenico: i principi di
    unificazione formale.
\end{snippet}

\begin{snippetdefinition}{isomorfismo-gestalt-definition}{Isomorfismo}
    L'\emph{isomorfismo} è il principio secondo cui l'organizzazione degli
    eventi psicologici è riconducibile alle proprietà strutturali degli eventi
    neurofisiologici corrispondenti.
\end{snippetdefinition}

\end{document}
